% Source: https://artofproblemsolving.com/community/c3h2796294_2017_tx_mmc__texas_mathworks_math_contest
% Stable identity: aops-topic-2796294-q13
Decide first which vertices are chosen. Two adjacent vertices cannot both be chosen (their common side would contain two chosen points). Conversely, given a set of pairwise non-adjacent vertices: a side with a chosen endpoint must not use its midpoint, and a side with no chosen endpoint must use its midpoint. So each such vertex set gives exactly one inner polygon, and every polygon has at least $4$ chosen points.

For example, choosing vertices $V_0,V_2,V_5$ (numbered around the octagon) forces the midpoints of sides $V_3V_4$ and $V_6V_7$, giving the pentagon below.
\begin{asy}
usepackage("amsmath");
size(3cm);
defaultpen(fontsize(10pt));
for(int i=0;i<8;++i){pair P=dir(22.5+45*i),Q=dir(22.5+45*(i+1)); draw(P--Q); dot(P); dot((P+Q)/2,linewidth(3));}
pair V0=dir(22.5),V2=dir(22.5+90),V5=dir(22.5+225);
pair m3=(dir(22.5+135)+dir(22.5+180))/2,m6=(dir(22.5+270)+dir(22.5+315))/2;
draw(V0--V2--m3--V5--m6--cycle,dashed+linewidth(1));dot(V0,linewidth(6));dot(V2,linewidth(6));dot(V5,linewidth(6));dot(m3,linewidth(6));dot(m6,linewidth(6));
label("$V_0$",V0,NE);label("$V_2$",V2,NW);label("$V_5$",V5,S);
\end{asy}

The number of ways to choose non-adjacent vertices of an $n$-cycle is the Lucas number $L_n$, which satisfies $L_n=L_{n-1}+L_{n-2}$ with $L_1=1$, $L_2=3$ (a path argument: condition on whether a fixed vertex is chosen). Thus $L_3=4$, $L_4=7$, $L_5=11$, $L_6=18$, $L_7=29$, $L_8=47$.

There are $\boxed{47}$ inner polygons.
